\documentclass[12pt,reqno]{amsart}
\usepackage{math-review}
\usepackage{needspace}
% Shared conventions only; this is not a manuscript.
% Smooth connected Riemannian manifolds, n >= 2, without boundary unless specified.
% R(X,Y)Z = nabla_X nabla_Y Z - nabla_Y nabla_X Z - nabla_[X,Y] Z.
% Ric is unnormalized; Ric >= (n-1) k g. Delta = div grad.
% D_k = pi/sqrt(k) for k > 0, infinity otherwise.
% s_k'' + k s_k = 0, s_k(0)=0, s_k'(0)=1; c_k=s_k'/s_k on (0,D_k).
% A is a radial Jacobi tensor; S=A'A^{-1} is the shape operator.
% m=tr S; normalized mean curvature a=m/(n-1).
% Busemann convention: b_gamma(x)=lim_{t->infinity}(t-d(x,gamma(t))).
% Canonical fragments use standard amsmath/amssymb commands, no private macros.

\title{Seven comparison questions in Riemannian geometry}
\author{}
\date{}
\begin{document}
\raggedbottom
\maketitle

Curvature controls how geodesics separate. Directional estimates give Rauch and Hessian comparison; their trace gives Laplacian and volume comparison. Global conclusions use endpoint fields, covering spaces or triangle gluing.

We assume smooth manifolds, linear algebra, calculus and elementary ODEs. Each of the seven theorems comes with its operative argument and explicit imported inputs.

\section{How fast can geodesics separate?}
\label{sec:rauch}

\subsection*{The objects and the comparison result}
A metric assigns inner products to tangent spaces. Curve length integrates speed, and coordinate volume has density \(\sqrt{\det(g_{ij})}\).
We use the torsion-free, metric-compatible Levi--Civita connection \(\nabla\).

Parallel transport solves \(D_tE=0\) and preserves inner products. In a parallel frame, covariant derivatives are component derivatives. A geodesic has \(D_t\dot\gamma=0\), and \(\exp_p(tv)\) follows initial velocity \(v\).

Our curvature convention is
\[
R(X,Y)Z=\nabla_X\nabla_YZ-\nabla_Y\nabla_XZ-\nabla_{[X,Y]}Z.
\]
For unit \(T\), the symmetric radial operator is \(\mathcal R_T X=R(X,T)T\) on \(T^\perp\). Its quadratic form is \(K(T\wedge X)|X|^2\), and its trace is \(\operatorname{Ric}(T,T)\).

A normal Jacobi field satisfies \(J\perp\dot\gamma\) and
\[
D_t^2J+R(J,\dot\gamma)\dot\gamma=0.
\]
It describes the first-order separation of a family of geodesics. The following constant-curvature comparison specifies both its scale and its permitted interval \cite{S-Eschenburg-1994}.

\begin{samepage}
\begin{theorem}[Rauch]
\label{thm:rauch}
Let \((M^n,g)\) be a smooth Riemannian manifold, \(n\ge2\), and \(\gamma:[0,T]\to M\) a unit speed geodesic. Let \(J\) be a normal Jacobi field with \(J(0)=0\) and \(|D_tJ(0)|=a>0\). For \(k\in\mathbb R\) let \(s\) solve
\(s''+ks=0,\ s(0)=0,\ s'(0)=1\), and set \(D_k=\pi/\sqrt{k}\) if \(k>0\) and \(D_k=+\infty\) otherwise.
If every two-plane containing \(\dot\gamma(t)\) has sectional curvature at most \(k\) for all \(t\in[0,T]\), then
\[
|J(t)|\ge a\,s(t)\qquad(0<t<\min\{T,D_k\}).
\]
If instead all those sectional curvatures are at least \(k\), and \(\gamma\) has no point conjugate to \(\gamma(0)\) in \((0,T]\), then
\[
|J(t)|\le a\,s(t)\qquad(0<t<\min\{T,D_k\}).
\]
\end{theorem}
\end{samepage}

Write this scalar solution as \(s_k\), and put \(c_k=s_k'/s_k\) wherever \(s_k>0\). Explicitly,
\[
s_k(t)=\begin{cases}
\sin(\sqrt{k}t)/\sqrt{k},&k>0,\\
t,&k=0,\\
\sinh(\sqrt{-k}t)/\sqrt{-k},&k<0.
\end{cases}
\qquad 0<t<D_k.
\]
In constant curvature, a unit initial normal derivative produces the field \(s_kE\). The positive interval is essential: the spherical model vanishes at \(D_k\).

\subsection*{From a variation to relative expansion}
Fix \(|v|=1\) and vary it with initial derivative \(w\perp v\). Differentiating \(\Phi(t,v)=\exp_p(tv)\) gives a Jacobi field with initial data \(J_w(0)=0\), \(J_w'(0)=w\). Hence
\[
J_w(t)=A(t)w,\qquad A''+\mathcal RA=0,\qquad A(0)=0,\quad A'(0)=I.
\]
The angular differential is \(d\Phi(w)=A(t)w\). A vector \(w\) at \(tv\in T_pM\), however, gives \(d(\exp_p)_{tv}(w)=A(t)w/t\); differentiating \(tv(s)\) accounts for the factor \(t\).

Where \(A\) is invertible, a current separation \(X=Aw\) has velocity \(A'w=A'A^{-1}X\). Thus relative expansion is
\[
S=A'A^{-1},\qquad S'+S^2+\mathcal R=0,\qquad S=t^{-1}I+O(t).
\]
The Taylor estimate follows from \(A=tI+O(t^3)\). The Wronskian \(A^TA'-(A')^TA\) is constant and zero, so multiplication by \(A^{-T}\) and \(A^{-1}\) proves \(S=S^T\).

\subsection*{The imported order and its actual use}
We import singular matrix Riccati comparison: if \(\mathcal R_1\ge\mathcal R_2\), their radial symmetric solutions \(S_i=t^{-1}I+O(t)\) obey \(S_1\le S_2\) on their common existence interval \cite[Remark 3.2]{S-Eschenburg-1994}.

Comparing with \(S_k=c_kI\) gives \(S\le c_kI\) under the sectional lower bound. For \(J=Aw\ne0\), the decisive operation is
\[
\frac d{dt}\log\frac{|J|}{|w|s_k}
=\frac{\langle SJ,J\rangle}{|J|^2}-c_k\le0,
\qquad \lim_{t\downarrow0}\frac{|J|}{|w|s_k}=1.
\]
Integration proves the upper length bound. Reversing the sectional hypothesis reverses this calculation. Its positive lower bound prevents \(Aw\) from first vanishing before \(D_k\), which explains the different conjugate-point requirements.

For an admissible curve of initial directions \(v(u)\), the same estimate bounds actual curve speed:
\[
\left|\frac d{du}\exp_p(tv(u))\right|
\le s_k(t)|v'(u)|.
\]
Integration bounds its image length. The rays must all satisfy the sectional lower bound and remain before their first conjugate point. This local length conclusion alone is not a global triangle theorem.

\section{What does this say about the distance Hessian?}
\label{sec:hessian}

For \(r=d(p,\cdot)\), the cut time \(\tau_p(v)\) is the last minimizing time of \(\exp_p(tv)\). The cut locus consists of its finite endpoints. A conjugate time instead means that a nonzero Jacobi field vanishes at both ends.

We import the regular-domain facts: no conjugate point occurs before cut; away from \(p\) and its cut locus, \(r\) is smooth with a unique minimizing ray \cite[Section 5]{S-Eschenburg-1994}.

The gradient represents \(dr\), and \(\nabla^2r(X,Y)=\langle\nabla_X\nabla r,Y\rangle\). This is a bilinear form, whereas \(S\) acts on normal vectors.

\begin{samepage}
\begin{theorem}[Hessian comparison]
\label{thm:hessian}
Let \((M^n,g)\) be a connected complete smooth Riemannian manifold, \(n\ge2\), \(p\in M\), and \(x\notin\{p\}\cup\operatorname{Cut}(p)\). Set \(r=d_g(p,\cdot)\) and \(\ell=r(x)>0\). Fix \(k\in\mathbb R\) and assume \(\ell<D_k\), where \(D_k=\pi/\sqrt{k}\) for \(k>0\) and \(+\infty\) otherwise. Let \(s\) solve \(s''+ks=0,s(0)=0,s'(0)=1\).
If the sectional curvatures of every plane containing the tangent to the unit speed minimizing segment from \(p\) to \(x\) are at least \(k\), then, for all \(X\in T_xM\),
\[
\nabla^2r_x(X,X)\le \frac{s'(\ell)}{s(\ell)}\bigl(|X|^2-dr_x(X)^2\bigr).
\]
If those radial sectional curvatures are instead at most \(k\), the inequality is reversed.
\end{theorem}
\end{samepage}

Gauss lemma identifies \(\nabla r=T=\dot\gamma\) on the uncut domain. For a radial variation, torsion freeness gives \(\nabla_JT=\nabla_TJ\). A specified current normal vector \(X\) therefore satisfies
\[
w=A^{-1}X,\qquad \nabla_X\nabla r=A'w=SX.
\]
The radial component is zero because \(\nabla_TT=0\). Decompose \(X=X^\perp+dr(X)T\); the matrix order from the first unit now gives
\[
\nabla^2r(X,X)=\langle SX^\perp,X^\perp\rangle
\le c_k(r)\bigl(|X|^2-dr(X)^2\bigr).
\]
This identifies exactly which operator has been estimated. At \(k=0\), the normal bound is \(|X|^2/r\); in Euclidean space it is an equality.

The smooth-domain restriction cannot be omitted. On a flat cylinder of circle radius \(a\),
\[
r([\theta],z)=\min_{m\in\mathbb Z}
\sqrt{(\theta+2\pi am)^2+z^2}.
\]
At \((\theta,z)=(\pi a,0)\), the angular one-sided derivatives are \(+1\) and \(-1\). There is no smooth distance Hessian there, although \(J=tE\) and no conjugate points occur.

\Needspace{14\baselineskip}
\section{What survives if only Ricci is bounded?}
\label{sec:laplacian}

Ricci bounds the sum of radial sectional curvatures. It need not bound each direction. The appropriate output is therefore the trace \(\Delta r=\operatorname{tr}_g\nabla^2r\), with \(\Delta=\operatorname{div}\nabla\).

\begin{samepage}
\begin{theorem}[Laplacian comparison]
\label{thm:laplacian}
Let \((M^n,g)\) be a connected complete smooth Riemannian manifold, \(n\ge2\), with \(\operatorname{Ric}\ge(n-1)k g\), \(k\in\mathbb R\). For every \(p\in M\), set \(r(x)=d_g(p,x)\). Let \(s\) be the solution of \(s''+ks=0,s(0)=0,s'(0)=1\), and let \(D_k=\pi/\sqrt{k}\) for \(k>0\), \(+\infty\) otherwise. With \(\Delta=\operatorname{div}\nabla\),
\[
\Delta r(x)\le(n-1)\frac{s'(r(x))}{s(r(x))}
\]
for every \(x\notin\{p\}\cup\operatorname{Cut}(p)\) with \(0<r(x)<D_k\).
\end{theorem}
\end{samepage}

Put \(m=\operatorname{tr}S\). Symmetry of \(S\) and Cauchy--Schwarz for its \(n-1\) eigenvalues give
\[
m'+\operatorname{tr}(S^2)+\operatorname{Ric}(T,T)=0,
\qquad \operatorname{tr}(S^2)\ge\frac{m^2}{n-1}.
\]
Consequently \(a=m/(n-1)\) obeys \(a'+a^2+k\le0\). Both \(a\) and \(c_k\) equal \(t^{-1}+O(t)\), so their singular initial behavior needs attention.

Set \(u=a-c_k\) and \(\mu=\exp\int_{t_*}^t(a+c_k)\), with \(t_*\) interior. Then
\[
(\mu u)'\le0,\qquad \mu\sim Ct^2,\qquad u=O(t),
\qquad \mu(t)u(t)\le\lim_{\varepsilon\downarrow0}\mu(\varepsilon)u(\varepsilon)=0.
\]
Thus \(m\le(n-1)c_k\). This is a scalar proof and does not import the singular matrix comparison from the first unit.

A Ricci upper bound cannot reverse the argument. Writing \(Q=\operatorname{tr}(S^2)-m^2/(n-1)\ge0\) leaves
\[
m'=-\frac{m^2}{n-1}-Q-\operatorname{Ric}(T,T).
\]
An upper bound for Ricci gives no lower control on \(-Q\). This identifies a failed inference; it does not by itself disprove every possible theorem involving a Ricci upper bound.

\section{How does one ray control an entire ball?}
\label{sec:volume}

Let \(\omega_{n-1}\) be the area of the unit sphere in \(\mathbb R^n\). The comparison ball volume is
\[
V_k^n(r)=\omega_{n-1}\int_0^{\min(r,D_k)}s_k(t)^{n-1}\,dt.
\]
The spherical model saturates at \(D_k\); its density is used only before that endpoint.

\begin{samepage}
\begin{theorem}[Bishop--Gromov]
\label{thm:volume}
Let \((M^n,g)\) be a connected complete smooth Riemannian manifold, \(n\ge2\), with \(\operatorname{Ric}\ge(n-1)k g\) for a fixed \(k\in\mathbb R\). Let \(V_k^n(r)\) denote the volume of a radius-\(r\) ball in the complete simply connected space form of dimension \(n\) and sectional curvature \(k\); for \(k>0\) it is constant for \(r\ge\pi/\sqrt{k}\).
For every \(p\in M\) and all \(0<r<R<\infty\),
\[
\frac{\operatorname{Vol}_g B(p,R)}{V_k^n(R)}
\le \frac{\operatorname{Vol}_g B(p,r)}{V_k^n(r)}\le1,
\qquad
\lim_{r\downarrow0}\frac{\operatorname{Vol}_g B(p,r)}{V_k^n(r)}=1.
\]
Equivalently,
\[
\frac{\operatorname{Vol}_g B(p,R)}{\operatorname{Vol}_g B(p,r)}\le\frac{V_k^n(R)}{V_k^n(r)}.
\]
\end{theorem}
\end{samepage}

\subsection*{Density, counting and averaging}
The radial differential is unit and orthogonal to the angular differential by Gauss lemma. Thus its metric Gram matrix and positive uncut volume density are
\[
\begin{pmatrix}1&0\\0&A^TA\end{pmatrix},
\qquad j=\sqrt{\det(A^TA)}=\det A>0.
\]
The last sign follows from \(A=tI+O(t^3)\) and absence of a zero before cut. Jacobi's determinant formula now converts the preceding trace estimate into
\[
(\log j)'=m,\qquad
\left(\log\frac{j}{s_k^{n-1}}\right)'\le0,
\qquad \frac{j}{s_k^{n-1}}\longrightarrow1.
\]
The model density removes the model's own expansion. This raywise estimate still needs a way to count ball points exactly once.

We import the cut-domain theorem: \(\exp_p\) maps \(\{tv:0<t<\tau_p(v)\}\) diffeomorphically onto the complement of \(p\) and its cut locus; the omitted set has volume zero. The measurable cut time then gives
\[
\operatorname{Vol}B(p,r)=\int_{S_pM}\int_0^{\min(r,\tau_p(v))}j(t,v)\,dt\,dv.
\]
Define \(h=j\) before cut and \(h=0\) at and after cut. For \(t<D_k\), the quotient \(h/s_k^{n-1}\) can only jump downward. Its angular average
\[
q(t)=\frac1{\omega_{n-1}}\int_{S_pM}\frac{h(t,v)}{s_k(t)^{n-1}}\,dv
\]
is nonincreasing and lies in \([0,1]\). Domination by one on the finite-measure sphere yields \(q(0+)=1\).

For \(r<D_k\), the ball ratio is a weighted average:
\[
F(r)=\frac{\int_0^r q(t)w(t)\,dt}{W(r)},
\qquad w=s_k^{n-1},\qquad W(r)=\int_0^r w(t)\,dt.
\]
Since \(q(t)\ge q(r)\) for \(t<r\), we have \(F(r)\ge q(r)\). On compact intervals away from zero, this quotient is absolutely continuous and
\[
F'(r)=\frac{w(r)}{W(r)}\bigl(q(r)-F(r)\bigr)\le0
\quad\text{almost everywhere}.
\]
Integrating proves monotonicity at every radius. The limit of the weighted averages is one. If \(k\le0\), this already covers all finite radii.

If \(k>0\), the independent Myers argument in the next unit gives \(\operatorname{diam}M\le D_k\). Points at distance \(D_k\) are cut points, so form a null set; both ball volumes saturate there. Continuity closes the endpoint without dividing by zero.

For \(\operatorname{Ric}\ge0\), taking \(R=2r\) gives
\[
\operatorname{Vol}B(p,2r)\le2^n\operatorname{Vol}B(p,r).
\]
This relative estimate does not prevent collapse: on \(S^1_\varepsilon\times\mathbb R^{n-1}\), a radius-\(r\) ball has volume at most \(2\pi\varepsilon\) times the Euclidean \((n-1)\)-ball volume, which tends to zero.

\section{Why does positive Ricci force a diameter bound?}
\label{sec:myers}

\begin{samepage}
\begin{theorem}[Bonnet--Myers]
\label{thm:myers}
Let \((M^n,g)\) be a connected complete smooth Riemannian manifold without boundary, \(n\ge2\). Suppose that there is a constant \(k>0\) such that
\[
\operatorname{Ric}_x(v,v)\ge(n-1)k|v|_g^2\quad\text{for every }x\in M\text{ and }v\in T_xM.
\]
Then \(M\) is compact,
\[
\operatorname{diam}_g(M)\le\pi/\sqrt{k},
\]
and \(\pi_1(M)\) is finite.
\end{theorem}
\end{samepage}

The input is fixed-end second variation: along a minimizing unit-speed segment \(\gamma:[0,L]\to M\), every normal field with \(V(0)=V(L)=0\) satisfies
\[
I(V,V)=\int_0^L\bigl(|D_tV|^2-\langle R(V,\dot\gamma)\dot\gamma,V\rangle\bigr)dt\ge0.
\]
This is an imported variational formula \cite[Section 7.1]{S-Bishop-2013v2}. Its test fields and the following integral are the operative proof.

Choose a parallel orthonormal normal frame and put \(V_i=\sin(\pi t/L)E_i\). These fields meet both endpoint conditions. Summing the curvature terms produces Ricci, so
\[
\begin{aligned}
\sum_{i=1}^{n-1}I(V_i,V_i)
&\le(n-1)\int_0^L\left[\frac{\pi^2}{L^2}\cos^2\frac{\pi t}{L}
-k\sin^2\frac{\pi t}{L}\right]dt\\
&=\frac{(n-1)L}{2}\left(\frac{\pi^2}{L^2}-k\right).
\end{aligned}
\]
For \(L>D_k\) the sum is negative, contradicting nonnegativity of every term. Hopf--Rinow supplies a minimizing segment between each pair and compactness of closed bounded subsets; hence the diameter bound makes \(M\) compact.

The universal cover with the lifted metric is complete: project, extend and lift each geodesic. It has the same Ricci lower bound, so it too is compact. Its fiber is closed and discrete, hence finite, and identifies with \(\pi_1(M)\).

These covering facts are imported topology. The open round hemisphere shows why completeness matters for compactness. A family of spheres of radius \(R\to\infty\) has positive Ricci but no uniform diameter bound without one common positive \(k\).

\section{When are global exponential coordinates possible?}
\label{sec:hadamard}

\begin{samepage}
\begin{theorem}[Cartan--Hadamard]
\label{thm:hadamard}
Let \((M^n,g)\) be a connected complete smooth Riemannian manifold with \(\sec\le0\). For every \(p\in M\), the exponential map
\[
\exp_p:T_pM\longrightarrow M
\]
is a covering map. If \(M\) is simply connected, \(\exp_p\) is a global diffeomorphism for every \(p\), \(M\) is diffeomorphic to \(\mathbb R^n\), and every pair of points is joined by a unique minimizing geodesic.
\end{theorem}
\end{samepage}

For a normal Jacobi field with \(J(0)=0\), metric compatibility and the Jacobi equation give
\[
\frac12(|J|^2)''=|J'|^2-\langle R(J,\dot\gamma)\dot\gamma,J\rangle\ge0.
\]
The sectional sign is used here. If \(J'(0)\ne0\), a second zero would force this nonnegative convex function to vanish between its zeros, contradicting \(J(t)=tJ'(0)+O(t^3)\).

Thus the angular tensor has no kernel. For \(w\perp v\), Gauss lemma also retains the radial direction:
\[
|d(\exp_p)_{tv}(av+w)|^2=a^2+t^{-2}|A(t)w|^2.
\]
The differential is invertible, including at zero. Completeness makes \(\exp_p\) defined on the whole tangent space, where \(h=\exp_p^*g\) is now a positive-definite metric.

Every \(h\)-geodesic from zero is a radial line: its image is the geodesic with the same initial velocity, and local-isometry preservation and ODE uniqueness identify the lift. All these lines are entire, so the one-point Hopf--Rinow criterion makes \((T_pM,h)\) complete.

For clarity, the complete-local-isometry covering input has the following mechanism. Over a small normal ball about \(y\), every \(x\in\exp_p^{-1}(y)\) supplies the inverse section
\[
z\longmapsto\exp_x^h\left((d\exp_p|_x)^{-1}(\exp_y^g)^{-1}(z)\right).
\]
Completeness defines the lifted radial segments. Reversing them and using geodesic uniqueness proves that these sections have disjoint images and exhaust the preimage of the ball. Thus they are full covering sheets.

A connected cover over a simply connected base has one sheet. Applying the result at every base point gives unique minimizing geodesics. On the flat cylinder the covering has many sheets, locating exactly where the extra topological hypothesis is needed.

\Needspace{24\baselineskip}
\section{How does local comparison become a triangle theorem?}
\label{sec:toponogov}

An SSS comparison fixes three side lengths. A same-angle, or SAS, hinge fixes two lengths and their included angle. These different fixed data explain the apparently opposite inequalities in the following theorem \cite{S-Lang-2026,S-Toponogov-1959}.

\begin{samepage}
\begin{theorem}[Toponogov]
\label{thm:toponogov}
Let \((M^n,g)\) be a connected complete smooth Riemannian manifold without boundary, \(n\ge2\), with \(\sec_M\ge k\). Put \(D_k=\pi/\sqrt{k}\) for \(k>0\) and \(D_k=+\infty\) otherwise. Choose distinct points \(p,x,y\) and minimizing segments \(px,py,xy\). Suppose their side lengths satisfy strict triangle inequalities and
\[
d(p,x)+d(p,y)+d(x,y)<2D_k.
\]
Let \((\bar p,\bar x,\bar y)\) be the triangle with the same three side lengths in the complete simply connected surface \(M_k^2\) of sectional curvature \(k\). Then each angle of the chosen triangle in \(M\) is at least the corresponding model angle.
If \((\hat p,\hat x,\hat y)\) is instead a model hinge with
\[
d(\hat p,\hat x)=d(p,x),\quad d(\hat p,\hat y)=d(p,y),\quad
\angle\hat x\hat p\hat y=\angle xpy,
\]
then \(d_M(x,y)\le d_{M_k^2}(\hat x,\hat y)\).
Moreover, for all \(u\in px\), \(v\in py\) and matching points \(\bar u\in\bar p\bar x\), \(\bar v\in\bar p\bar y\) such that
\(d(p,u)=d(\bar p,\bar u)\) and \(d(p,v)=d(\bar p,\bar v)\),
\[
d_M(u,v)\ge d_{M_k^2}(\bar u,\bar v).
\]
\end{theorem}
\end{samepage}

\subsection*{The local map and its permitted curve}
Choose a common normal radius \(R<D_k\) and a linear tangent-space isometry \(I\) from the \(n\)-dimensional model to \(M\). The map
\[
F=\exp_p\circ I\circ(\exp_o^k)^{-1}:B_k(o,R)\longrightarrow B_M(p,R)
\]
preserves radial speed and contracts angular Jacobi fields by Rauch. Gauss lemma therefore gives \(|dF(U)|\le|U|\), including mixed radial and angular vectors.

Take endpoints with radial lengths \(a,b<R/2\), and a minimizing model chord of length \(L\le a+b\). At arclength \(s\) along that chord,
\[
d_k(o,\bar\sigma(s))\le\min(a+s,b+L-s)
\le\frac{a+b+L}{2}\le a+b<R.
\]
The chord remains inside the map's domain. Integrating speed proves the local SAS distance inequality; the strictly increasing model opposite-side function of the included angle gives local SSS angle comparison.

\subsection*{The global input and the shrinking contradiction}
We import Lang's Proposition 5.14. Its hypotheses are balanced segments and existence of minimizing segments for pairs at distance below \(D_k\).

Balanced means that angles to the two opposite directions on a segment sum to \(\pi\). This holds for smooth tangent angles; Hopf--Rinow supplies the required segments.

For a hinge \(H_p(x,y)\) with nonzero sides and perimeter \(P<2D_k\), if every hinge of perimeter below \(4P/5\), with a radial-side endpoint in \(px\cup\{y\}\), has SSS angle comparison, then \(H_p(x,y)\) does too.

Auxiliary degenerate triples use their continuous comparison angles. The principal theorem keeps strict triangle inequalities. This gluing statement is a technical input, not a consequence of the local map alone.

If comparison failed, the contrapositive would produce violating hinges with
\[
P_j<(4/5)^jP_0,\qquad
d(p_j,p_{j-1})\le P_j+P_{j-1}
<\tfrac95(4/5)^{j-1}P_0.
\]
The shared endpoint gives the second estimate. Summing yields \(d(p_m,p_n)\le9P_0(4/5)^n\), so completeness supplies a limit \(q\); the endpoints approach it as well.

Local comparison must hold uniformly at the moving vertices. The inverse function theorem applied to \((p,\xi)\mapsto(p,\exp_p\xi)\), followed by Gauss lemma, supplies one normal radius for all \(p\) near \(q\). The small violating hinges eventually lie there, a contradiction.

\subsection*{Why the chord conclusion needs more than one angle}
We also import Lang's Lemma 5.9: angle comparison for all auxiliary hinges connecting points of \(px\cup py\) implies SSS chord comparison.

This technical conversion still requires a perimeter check before the global angle result can be applied to its auxiliary hinges.

Let \(u_1,u_2\in px\) have distances \(r_1\le r_2\) from \(p\), and \(v\in py\) have distance \(s\). With \(a=d(p,x)\), \(b=d(p,y)\), \(c=d(x,y)\),
\[
P(u_1,u_2,v)\le(r_2-r_1)+(r_1+s)+(a-r_2+c+b-s)=a+b+c.
\]
Every triple has two points on one side, so the global angle theorem covers all required auxiliary hinges. Completeness and balanced sides supply the other inputs. This proves the claimed chord comparison with its exact strict perimeter boundary.

For \(k=0\), the SAS conclusion reads \(d(x,y)^2\le a^2+b^2-2ab\cos\theta\). It cannot simply be reversed under \(\sec\le0\): on \(\mathbb R^2/\mathbb Z^2\), lifts \((0.4,0.1)\), \((-0.4,0.1)\) give model distance \(0.8\) but torus distance \(0.2\).

The seven questions thus share radial variations but consume different global inputs. General rigidity, the spherical perimeter endpoint, weighted curvature and synthetic extensions require additional statements and are outside this account.

\bibliographystyle{amsplain}
\bibliography{../generated/references}
\end{document}
